% Worked Examples: Concept 2.5.2 - Determining the Transfer Function
\documentclass[11pt,a4paper]{article}
\usepackage{worked-examples-template}

\title{\textbf{Worked Examples}\\[0.5em]
\large Concept 2.5.2: Determining the Transfer Function\\[0.3em]
\normalsize \coursesubtitle{} -- \coursetitle}
\author{Assoc.\ Prof.\ William Robertson \and Dr Sean McGowan}
\date{\today}

\begin{document}

\maketitle
\tableofcontents
\newpage

\section{Concept 2.5.2: Transfer Functions}

\subsection{Example 1: Deriving Transfer Function from State-Space Model}

\paragraph{Problem:} Find the transfer function $G(s) = Y(s)/U(s)$ for the state-space system:
\begin{align}
\frac{dx}{dt} = \begin{bmatrix} -2 & 1 \\ 0 & -3 \end{bmatrix}x + \begin{bmatrix} 1 \\ 1 \end{bmatrix}u, \quad y = \begin{bmatrix} 1 & 2 \end{bmatrix}x
\end{align}

\paragraph{Solution:}

The transfer function from state-space is given by:
\begin{align}
G(s) = C(sI - A)^{-1}B + D
\end{align}

where $D = 0$ (no direct feedthrough).

First, compute $sI - A$:
\begin{align}
sI - A = \begin{bmatrix} s & 0 \\ 0 & s \end{bmatrix} - \begin{bmatrix} -2 & 1 \\ 0 & -3 \end{bmatrix} = \begin{bmatrix} s+2 & -1 \\ 0 & s+3 \end{bmatrix}
\end{align}

Find the inverse using the adjugate method. For a $2 \times 2$ matrix, the adjugate (also called adjoint) is the transpose of the cofactor matrix:
\begin{align}
(sI-A)^{-1} = \frac{1}{\det(sI-A)}\text{adj}(sI-A)
\end{align}

Determinant:
\begin{align}
\det(sI-A) = (s+2)(s+3) = s^2 + 5s + 6
\end{align}

Adjugate (swap diagonal, negate off-diagonal):
\begin{align}
\text{adj}(sI-A) = \begin{bmatrix} s+3 & 1 \\ 0 & s+2 \end{bmatrix}
\end{align}

Therefore:
\begin{align}
(sI-A)^{-1} = \frac{1}{s^2+5s+6}\begin{bmatrix} s+3 & 1 \\ 0 & s+2 \end{bmatrix}
\end{align}

Compute $(sI-A)^{-1}B$:
\begin{align}
(sI-A)^{-1}B &= \frac{1}{s^2+5s+6}\begin{bmatrix} s+3 & 1 \\ 0 & s+2 \end{bmatrix}\begin{bmatrix} 1 \\ 1 \end{bmatrix} \\
&= \frac{1}{s^2+5s+6}\begin{bmatrix} s+4 \\ s+2 \end{bmatrix}
\end{align}

Finally, multiply by $C$:
\begin{align}
G(s) &= C(sI-A)^{-1}B \\
&= \begin{bmatrix} 1 & 2 \end{bmatrix}\frac{1}{s^2+5s+6}\begin{bmatrix} s+4 \\ s+2 \end{bmatrix} \\
&= \frac{(s+4) + 2(s+2)}{s^2+5s+6} \\
&= \frac{s+4+2s+4}{s^2+5s+6} \\
&= \frac{3s+8}{s^2+5s+6} \\
&= \frac{3s+8}{(s+2)(s+3)}
\end{align}

\textbf{Result:} The transfer function is:
\begin{align}
G(s) = \frac{3s+8}{(s+2)(s+3)}
\end{align}

The poles are at $s = -2$ and $s = -3$ (eigenvalues of $A$).
The zero is at $s = -8/3 \approx -2.67$.
\end{document}
